\section{决赛运行时设计}

\subsection{阶段感知的映射建议}

一次性对整个模型设置固定页建议并不可靠。Prefill 通常按层扫描较大范围，连续预读可以摊薄缺页与存储请求开销；Decode 只激活少量专家，但在内存远小于模型且存储随机访问较弱时，彻底关闭预读反而会把大块顺序 I/O 分裂为同步小读。SLIM-ARC 因而把页建议作为设备策略：加载与 Prefill 使用 \texttt{SEQUENTIAL}；Decode 在随机按需和保留预读之间由设备配置选择。该机制只改变内核 readahead 提示，不改变计算图或 Router 结果。

\subsection{层感知预取调度}

层级预取只覆盖本轮计算图及其有限前瞻窗口。设 $l_{\min}$ 和 $l_{\max}$ 为计算图中出现的最小、最大 Transformer 层，$w$ 为随 Prefill/Decode 阶段变化的前瞻层数，$T_l$ 为第 $l$ 层已注册的 mmap 权重张量集合。算法\ref{alg:layer-prefetch}给出调度路径；循环上界已经截断到 $l_{\max}$，因此不再保留旧稿中永远不可满足的 $l>l_{\max}$ 分支。

\begin{algorithm}[H]
  \caption{有界层感知预取}
  \label{alg:layer-prefetch}
  \renewcommand{\algorithmicrequire}{\textbf{输入:}}
  \renewcommand{\algorithmicensure}{\textbf{输出:}}
  \begin{algorithmic}[1]
    \REQUIRE 计算图 $G$，阶段相关窗口 $w$
    \ENSURE 对有限后续层发出 best-effort 页建议
    \STATE $(l_{\min},l_{\max}) \gets \mathrm{layer\_range}(G)$
    \STATE $l_{\mathrm{end}} \gets \min(l_{\max},l_{\min}+w)$
    \FOR{$l=l_{\min}$ \TO $l_{\mathrm{end}}$}
      \FORALL{$t \in T_l$}
        \STATE $\mathrm{madvise}(t.\mathrm{addr},t.\mathrm{size},\mathrm{WILLNEED})$
      \ENDFOR
    \ENDFOR
  \end{algorithmic}
\end{algorithm}

Router 观测不通过裸指针跨线程共享。每次预取分配单调递增的 generation token，并把层号、专家集合和成功建议字节作为不可变记录发布；实际 Router 结果只结算同 token 的记录。计算失败、缺少 Router 节点或过滤后集合为空时显式取消 token，从而避免旧预测占满有界队列。

\subsection{统一预算：计算方式与生效范围}

调度器先从静态周期额度 $B_{\mathrm{static}}$ 出发；压力准入开启时，根据 cgroup 当前占用、上限和保留量得到有效额度 $B_{\mathrm{eff}}\le B_{\mathrm{static}}$。随后按阶段系数 $\boldsymbol{\alpha}_p=(\alpha_w,\alpha_{kv},\alpha_e)$ 计算
\begin{equation}
  (B_w,B_{kv},B_e)=B_{\mathrm{eff}}\boldsymbol{\alpha}_p,
  \qquad \alpha_w+\alpha_{kv}+\alpha_e=1.
\end{equation}
运行时统计可在有限范围内调高发生 stall 的路径，并缩减另外两路。当前实现采用启发式乘法调整，没有在每次调整后对三路比例做二次归一化；因此这些额度应解释为各路径的 admission 上限，而不是三路实际 I/O 字节和严格等于 $B_{\mathrm{eff}}$。表\ref{tab:budget-ratios-finals}列出当前初始系数。

\begin{table}[H]
  \centering
  \caption{统一预算的阶段初始比例（策略参数，不是测得带宽占比）}
  \label{tab:budget-ratios-finals}
  \begin{tabular}{lccc}
    \toprule
    阶段 & 权重 $\alpha_w$ & KV $\alpha_{kv}$ & 专家 $\alpha_e$ \\
    \midrule
    PREFILL\_SHORT & 60\% & 10\% & 30\% \\
    PREFILL\_LONG  & 50\% & 20\% & 30\% \\
    DECODE\_SHORT  & 70\% & 20\% & 10\% \\
    DECODE\_LONG   & 30\% & 60\% & 10\% \\
    MOE\_DECODE    & 20\% & 20\% & 60\% \\
    \bottomrule
  \end{tabular}
\end{table}

必须区分“计算出额度”和“严格执行额度”。当前权重预取通过 \texttt{set\_memory\_budget} 执行 $B_w$，专家路径在压力准入、驻留或专家预算开关启用时执行 $B_e$；决赛的 model-owned runtime 尚未绑定 KV manager，因此 $B_{kv}$ 只存在于策略计算中，没有进入正式实验路径。即使后续绑定 KV manager，现有接口也尚未把 $B_{kv}$ 作为逐周期硬 admission。因此本报告将其称为\textbf{统一预算策略与两路执行路径}，不宣称权重、KV、专家三路已经完成同等强度的带宽隔离。补齐 KV manager 绑定、字节级 admission 和比例归一化后，才能把表中 KV 比例解释为实际约束。

\subsection{错误专家页回收：安全优先的候选建议}

每轮 Router 观测产生已预取集合和实际选中集合。回收规划器保留前者中未出现在后者的 ID，并保持首次出现顺序。对于每个候选专家，它验证张量的专家数、可整除布局、ID 范围和无溢出的偏移计算；再把专家切片向内收缩到完整页面。一个专家通常跨越多个页面，边界页还可能与相邻切片共享，因此系统只对完全位于该专家切片内部的页面提出建议。零长度、跨越张量边界、非整除布局或非法 ID 都不产生页面建议，而是记录可审核的原因。这样，建议绝不覆盖选中专家，也不触及专家切片外的字节。

\begin{figure}[H]
\centering
\begin{tikzpicture}[node distance=9mm and 7mm, >=Latex, every node/.style={font=\small,align=center,draw,rounded corners,minimum height=9mm}]
  \node[fill=blue!9] (router) {Router 观测\\\textbf{事实}};
  \node[fill=orange!12,right=of router] (plan) {差集与内部页规划\\\textbf{设计}};
  \node[fill=green!10,right=of plan] (advice) {可选页建议\\\textbf{设计}};
  \node[fill=gray!14,below=of plan] (metrics) {命中、浪费、跳过、失败\\\textbf{事实记录}};
  \draw[->] (router) -- (plan);
  \draw[->] (plan) -- (advice);
  \draw[->] (advice) |- (metrics);
  \draw[->] (router) -- (metrics);
\end{tikzpicture}
\caption{安全错误专家页回收闭环（\textbf{设计图，非性能测量}）。事实记录由严格运行时指标输出；是否带来性能收益必须由第\ref{sec:protocol}节协议的正式结果判定。}
% chktex 24 is a false positive: this figure label belongs after its caption.
\label{fig:reclaim-loop} % chktex 24
\end{figure}

\subsection{压力/准确性导向的专家驻留}

驻留不是无条件缓存。运行时从 cgroup v2 读取 \texttt{memory.current} 与 \texttt{memory.max}，将压力归为 normal、high、critical 或 missing；无效、溢出或不一致读数走显式兼容回退，而不被当作剩余容量。策略同时参考预测浪费的 EWMA 和饱和的专家热度：高压力时缩紧准入，错误率高时不为不可靠预测预留预算。压力阈值、恢复所需的连续样本和预算均为可测试的确定性规则，而非从最终结果反推的参数。

\subsection{不变量与设计取舍}

回收与驻留只发出建议，建议失败记为指标并回退，不使推理失败。其代价是保守：页边缘和不规则张量可能被跳过，错误预测也可能使驻留无收益。该取舍优先于以任何性能代价换取覆盖率；是否满足推广门槛留给受控实验决定。
